\section{One-step phase coercivity and the operator criterion}
\label{sec:operator-bridge}
The positive-measure estimate in Section~\ref{sec:periodic-coercivity}
uses an entire periodic tube. There is a stronger estimate if one first
locates a large one-step defect on the period, and only then thickens
that single state. It avoids the exponentially shrinking tube of the
whole period. We retain the constants $\rho_0,L,B_0$ of
Lemma~\ref{lem:uniform-period-tubes}, and set $r_0=\rho_0/(2L)$.

\begin{theorem}[One-step thickening of the periodic defect]
\label{thm:one-step-coercivity}
Fix $0<\alpha\le1$ and $1\le p<\infty$. Let $q$ be circle-valued and
measurable, with an $\alpha$-H\"older representative of constant $H\ge1$
on the $\rho_0$ balls around the states of the triple selected in
Theorem~\ref{thm:uniform-phase-separation}. For $b=|\beta|\ge1$ put
\[
 d_b=\frac1{4000000\pi M(b)},\qquad
 H_b=H(1+L^\alpha),\qquad
 r_b=\min\{r_0,d_b/(4bB_0),(d_b/(4H_b))^{1/\alpha}\}.
\]
For the actual one-step defect $\mathcal D_q$ defined in
Theorem~\ref{thm:positive-measure-coercivity},
\begin{equation}\label{eq:one-step-coercivity}
 \|\mathcal D_q\|_{L^p(\nu_R^*)}
       \ge\frac{d_b}{2}\left(\frac{r_b^2}{4c_*}\right)^{1/p}.
\end{equation}
If $H\le H_0b^\zeta$ with fixed $H_0\ge1$ and $\zeta\ge0$, this gives
\begin{equation}\label{eq:sharper-phase-power}
 \|\mathcal D_q\|_{L^p(\nu_R^*)}
 \ge C\,b^{-\lambda}(1+\log b)^{-\eta},\qquad
 \lambda=\frac2p\max\{1,\zeta/\alpha\},\quad
 \eta=1+\frac2{\alpha p}.
\end{equation}
The constants are uniform in the radius and all torus frequencies.
\end{theorem}
\begin{proof}
Telescoping the three exact periods as in
Corollary~\ref{cor:logarithmic-coercivity} shows that at some state $x_0$
of the selected triple the one-step defect has modulus at least $d_b$.
This is initially a statement about the specified continuous
representatives at finitely many states: the sum of the three induced
periods is at most $2M(b)$, while their product discrepancy is at least
$2/(4000000\pi)$.

On $B(x_0,r_0)$ the return has one fixed lattice label and one fixed
physical count, and $F_R^*$ maps that ball into the $\rho_0$ ball about
the successor of $x_0$. For $x$ in that ball,
\[
 |\mathcal D_q(x)-\mathcal D_q(x_0)|
 \le H_b|x-x_0|^\alpha+bB_0|x-x_0|.
\]
Each term on the right is at most $d_b/4$ on the ball of radius
$r_b$. On that ball the defect
is consequently at least $d_b/2$. Its actual invariant measure is
$r_b^2/(4c_*)$. Integration proves \eqref{eq:one-step-coercivity}.
Changes of representatives on null sets do not alter the integral.

Finally $M(b)\le3+\log b/\log(100/49)$ and
$H_b\le C b^\zeta$. The minimum defining $r_b$ is bounded
below by a positive constant times
$b^{-\max\{1,\zeta/\alpha\}}(1+\log b)^{-1/\alpha}$.
Substitution proves \eqref{eq:sharper-phase-power}. Unlike the
whole-period tube estimate, its exponent contains no $\log L$ term.
The regularity premise on $q$ remains essential.
\end{proof}

The next theorem specifies exactly what is required to turn this
physical coercivity into an operator resolvent. The hypothesis is a
reconstruction of a phase from an approximate spectral vector; it is
not silently identified with an eigenvalue equation.

\begin{theorem}[A quantitative phase-reconstruction criterion]
\label{thm:phase-resolvent-criterion}
For each $R,u,s,\beta$, let $\mathcal L_{R,u,s,\beta}$ be a bounded
operator on a Banach space $\mathcal B_R$. Fix the constants
$\alpha,p,H_0,\zeta$ of Theorem~\ref{thm:one-step-coercivity}, and let
$A\ge1$, $v\ge0$, $\theta>0$. Suppose that for every $b=|\beta|\ge1$,
$|z|=1$, and $h\in\mathcal B_R$ with $\|h\|=1$ and
$e=\|(z-\mathcal L_{R,u,s,\beta})h\|\le1$, one can construct a
circle-valued $q$ with the local H\"older bound $H_0b^\zeta$ and
\begin{equation}\label{eq:phase-reconstruction}
 \|\mathcal D_q\|_{L^p(\nu_R^*)}\le A b^v e^\theta,
 \qquad e^{ic}=z.
\end{equation}
Suppose also that $z-\mathcal L_{R,u,s,\beta}$ is Fredholm of index
zero. Then it is invertible and
\begin{equation}\label{eq:phase-resolvent}
 \|(z-\mathcal L_{R,u,s,\beta})^{-1}\|
 \le C b^{(\lambda+v)/\theta}
                (1+\log b)^{\eta/\theta},
\end{equation}
with $\lambda,\eta$ as in \eqref{eq:sharper-phase-power}.
\end{theorem}
\begin{proof}
Combining \eqref{eq:sharper-phase-power} and
\eqref{eq:phase-reconstruction} yields
\[
 e\ge (C_0/A)^{1/\theta}
       b^{-(\lambda+v)/\theta}(1+\log b)^{-\eta/\theta}
\]
when $e\le1$. Decreasing the positive constant to at most one makes
the same lower bound valid when $e>1$. Homogeneity gives a uniform
lower bound for $z-\mathcal L$ on every vector. It is therefore
injective and has closed range. Fredholm index zero makes its cokernel
zero as well, so it is surjective. The lower bound gives
\eqref{eq:phase-resolvent}. The Fredholm assumption cannot be omitted:
an injective operator bounded below need not be onto.
\end{proof}

For the unbounded induced twists, the phase reconstruction
\eqref{eq:phase-reconstruction}, including nonvanishing modulus and
quantitative regularity, is an unverified premise, not a result of
Lemma~\ref{lem:collision-spectral-input}. The latter treats a smooth
collision-space weight near zero. It provides the return-tail theorem
but is not a construction of $\mathcal B_R$ or of the raw high-frequency
operator. The criterion separates the necessary analytical construction
from the now explicit phase lower bound.
